%% English Abstract — appears immediately after the cover pages.
%% Within \begin{abstract} ... \end{abstract} environment provided by postech-ucs.cls.
%% Limit: 1,000 words (지침 1.나.③)

Recent advancements in the direct training of Spiking Neural Networks (SNNs) have demonstrated high-quality outputs even at early timesteps, paving the way for novel energy-efficient AI paradigms. However, the inherent non-linearity and temporal dependencies in SNNs introduce persistent challenges, such as temporal covariate shift (TCS) and unstable gradient flow with learnable neuron thresholds.

In this thesis, we enhance the stability of direct SNN training through two key innovations: \textbf{MP-Init (Membrane Potential Initialization)} and \textbf{TrSG (Threshold-Robust Surrogate Gradient)}. MP-Init addresses TCS by aligning the initial membrane potential with its stationary distribution, while TrSG stabilizes gradient flow with respect to the threshold voltage during training. We provide a rigorous mathematical analysis of both methods, prove the geometric convergence of the membrane potential to a stationary distribution, and characterize the threshold-induced gradient pathologies of existing surrogate-gradient families.

Extensive experiments validate our approach, achieving state-of-the-art accuracy on both static (CIFAR-10/100, ImageNet) and dynamic (DVS-CIFAR10) image datasets, while incurring negligible parameter and computational overhead. The methods further generalize cleanly to Transformer-based SNNs and to event-based object detection, indicating broad applicability beyond the classification setting.

The code for reproducing the results is available at: \href{https://github.com/kookhh0827/SNN-MP-Init-TRSG}{\textcolor{mypink}{https://github.com/kookhh0827/SNN-MP-Init-TRSG}}.
